\documentclass[11pt]{article}

\usepackage[margin=1in]{geometry}
\usepackage{amsmath, amssymb, amsthm}
\usepackage{xcolor}
\usepackage{booktabs}
\usepackage[colorlinks=true, linkcolor=blue!50!black, citecolor=blue!50!black, urlcolor=blue!50!black]{hyperref}

\newcommand{\F}{\mathbb{F}}
\newcommand{\K}{K}            % F_{2^64}, the data field
\newcommand{\E}{E}            % F_{2^128}, the proof field (degree-2 tower)
\newcommand{\eq}{\operatorname{eq}}
\newcommand{\cube}[1]{\{0,1\}^{#1}}
\newcommand{\gen}{\mathsf{g}}
\newcommand{\pkd}{\mathrm{pkd}}
\newcommand{\WP}[1]{\textbf{WP#1}}
\newtheorem{remark}{Remark}

\title{\bf Transition to 64 bits}
\author{}
\date{}

\begin{document}
\maketitle
\vspace{-3em}
\begin{center}
Moving leanVM-b's data to $\F_{2^{64}}$ inside the degree-2 tower $\F_{2^{128}}=\F_{2^{64}}[y]$:\\
64-bit machine words, commitments over $\F_{2^{64}}$, challenges in $\F_{2^{128}}$, one ring switching.
\end{center}

\tableofcontents

\section{Motivation and summary}\label{sec:summary}

The machine and its proof system currently live entirely in $\F_{2^{128}}$
(\texttt{misc/doc.tex}): every machine value is a 128-bit field element, and
the committed witness is $\F_{2^{128}}$-valued. But the VM's data (memory
words, addresses, registers, counters) does not need 128 bits; committing it
at that width doubles the witness for nothing.

We keep the 128-bit challenge field, and with it today's $\sim$120-bit
soundness, but split the roles that $\F_{2^{128}}$ plays today between the
two levels of a degree-2 \emph{tower}
\[
  \K=\F_{2^{64}},\qquad \E=\F_{2^{128}}=\K[y]/(y^2+y+c):
\]
\begin{itemize}
\item \textbf{Data and commitments live in $\K$.} Memory words, addresses,
  registers, operands, counts, opcodes, separators: everything the VM
  touches is a $\K$-element, and the committed witness is one dense
  $\K$-valued stack, Reed--Solomon encoded over $\K$ (8-byte symbols).
  Committed data \emph{halves} in size relative to today. Flock's Boolean
  witness stays in $\F_2$ and is packed $64$ bits per $\K$-word.
\item \textbf{Randomness lives in $\E$.} Every challenge ($\alpha,\gamma,
  \eta,\lambda$, zerocheck points, sumcheck and GKR rounds, PCS folds, and
  the ring-switching batch) is sampled from $\E$, so every interactive
  error term keeps its current form $c/2^{128}$: the security level does
  not move. Since $\K\subset\E$ is a basis-aligned subfield, $\K$-valued
  data combines with $\E$-valued challenges at a fraction of today's cost.
\item \textbf{One ring switching}~\cite{rsgen}, for flock only: the $\F_2$
  witness packs into $\K$ ($64=2^6$ bits per element) and its claims,
  opened over $\E$, reduce to one claim on the packed $\K$-polynomial. The
  VM columns need \emph{no} ring switching at all: they are committed
  directly, and their claims are $\E$-weighted sums over $\K$-data, which
  the PCS opens natively (\S\ref{sec:rs}).
\end{itemize}
Structurally the system stays exactly today's architecture (one stacked
witness, one claim pool, one batch-mixed Ligerito opening, one flock ring
switching): two type substitutions (data $\F_{2^{128}}\to\K$, randomness
$\F_{2^{128}}\to\E$) and \texttt{LOG\_PACKING} $7\to6$. This transition
buys witness size and prover speed, not security: a $\K$-product costs $1$
\texttt{PMULL} (versus $6$ for today's product) and the mixed
$\K\times\E$ product costs $2$.

\paragraph{Outlook.} When more than $128$ bits of soundness are needed,
swap $\E$ for the degree-3 extension $\F_{2^{192}}=\K[y]/(y^3+y+1)$
(implemented and benchmarked: type \texttt{F192}): commitments stay over
$\K$, only the challenge field grows, and every interactive error term
drops to $c/2^{192}$. This document targets the degree-2 step only.

\section{The field stack}\label{sec:fields}

\paragraph{The base field $\K$.}
$\K=\F_{2^{64}}=\F_2[x]/(x^{64}+x^4+x^3+x+1)$, the standard low-weight
irreducible pentanomial; reduction folds by $\mathtt{0x1B}$. One
multiplication is a single \texttt{PMULL} plus a fold. Crucially for the VM
(\S\ref{sec:vm-transport}), $x$ is \emph{primitive}: its order is exactly
$2^{64}-1$ (verified by checking $x^{(2^{64}-1)/q}\neq 1$ for every prime
$q\mid 2^{64}-1=3\cdot5\cdot17\cdot257\cdot641\cdot65537\cdot6700417$; a
unit test pins this).

\paragraph{The extension $\E$.}
$\E=\K[y]/(y^2+y+c)$ with $c=x^{61}$. By Artin--Schreier, $y^2+y+c$ is
irreducible over $\K$ exactly when $\mathrm{Tr}_{\K/\F_2}(c)=1$, and
$x^{61}$ is the least monomial of trace $1$ ($x^{63}$ is the only other
one; a unit test pins the choice). Elements are $e_0+e_1y$ with
$e_0,e_1\in\K$: two 64-bit lanes, the natural layout for packing pairs of
$\K$-values. This $\E$ is isomorphic to today's GHASH field but in a
different representation: transcripts and serialized proofs change bytes,
not security.

\paragraph{Embeddings and mixed arithmetic.}
$\F_2\subset\K\subset\E$; the embedding $\K\hookrightarrow\E$ is
$e\mapsto(e,0)$, free. The workhorse is the \emph{mixed} product
\[
  \K\times\E\to\E,\qquad k\cdot(e_0+e_1y)=ke_0+ke_1\,y,
\]
two base multiplications ($2$ \texttt{PMULL}): a third of the cost of
today's $\F_{2^{128}}$ product ($6$ \texttt{PMULL}). The full
$\E\times\E$ product is a 2-term Karatsuba ($3$ \texttt{PMULL} products,
the fold $y^2=y+c$ with the sparse constant $c=x^{61}$, and two base
reductions), on par with today's product; the exact schedule is fixed at
\WP1 and \texttt{field\_bench} arbitrates. In \texttt{PMULL} counts:
$\K\times\K$ costs $1$, $\K\times\E$ costs $2$, $\E\times\E$ about $6$.

\section{What lives where}\label{sec:vm-transport}

\subsection{Machine values move to $\K$}

Every VM quantity of \texttt{doc.tex} \S1 transports verbatim with
$\F_{2^{128}}$ replaced by $\K$ and $\gen=x\in\K$:
memory cells, the registers $\mathbf{pc},\mathbf{fp}$, operands, immediates,
read counts, domain separators, opcodes, all $\K$-elements; addresses and
counts remain powers of $\gen$, with $\operatorname{ord}(\gen)=2^{64}-1$.

\paragraph{Instance caps.}\label{sec:caps}
Several soundness arguments count objects against
$\operatorname{ord}(\gen)$. At $2^{128}$ they could be waved through,
since no instance approaches $2^{128}$ of anything; at $2^{64}$ that
hand-wave is not an argument, and we do not make it. Instead the protocol
fixes explicit \emph{caps} as public constants, and the verifier checks the
\emph{announced} instance against them before running any reduction,
rejecting on violation:
\[
  h\le 32,\qquad N\le 2^{32},\qquad \sum_{t}2^{\tau_t}\;\le\;2^{56}
\]
(memory size, program length, total announced rows across all tables). A
row makes at most $16$ reads (generously; the post-transition
\texttt{BLAKE3} row makes $13$: twelve memory, one bytecode), so every
\emph{admissible} instance performs
\[
  R\;\le\;16\cdot 2^{56}+2\cdot 2^{h}+2\cdot N\;<\;2^{61}
\]
read flushes in total: a bound that holds for every instance the
verifier accepts, not merely realistic ones. ($2^{56}$ rows is
$\sim\!10^{17}$, far beyond any provable trace, so the caps cost no
completeness; the exact constants may be revisited, what matters is that
they are checked and keep $R$ provably below $2^{64}-1$.)

The counting arguments are then theorems, conditional only on the cap
check:
\begin{itemize}
\item \emph{Memory soundness} (doc \S6.1): if the bus balances while some
  read of address $i$ returns $v\neq\mathbf{m}[i]$, the multiset of
  value-$v$ read counts is invariant under $\times\gen$, so nonempty it has
  size at least $\operatorname{ord}(\gen)=2^{64}-1>R$, contradicting the
  cap. Hence it is empty and every read returns the committed value.
\item \emph{Counts cannot wrap}: every count exponent is at most
  $R<2^{64}-1$, so the counts $\gen^{0},\gen^{1},\dots$ occurring in one
  address's history are pairwise distinct and the finalize count is
  uniquely determined.
\item \emph{Range checks in the exponent} (doc \S9.5): $e+f\equiv k-1
  \pmod{2^{64}-1}$ with $e,f<2^{h}\le2^{32}$ and $k\le 2^{16}$ forces
  $e+f<2^{33}<2^{64}-1$: the congruence is an equality over the integers,
  and $e\le k-1$ follows exactly as before.
\item \emph{Monomial index encoding}: $\gen^k=x^k$ is the $k$-th monomial for
  $k<64$ (previously $k<128$); the XMSS encoding check must budget indices
  accordingly.
\end{itemize}
When \texttt{doc.tex} is next revised, its ``automatic at any real
scale'' remarks (doc \S6.1) must be replaced by these caps.

\subsection{ISA deltas}

\begin{itemize}
\item \texttt{MUL\_NATIVE} becomes the $\K$-product (one \texttt{PMULL}
  natively; in-circuit constraint $v_C=v_A\cdot v_B$ is now a $\K$-relation,
  degree 2 as before).
\item \texttt{BLAKE3}: a 256-bit value now spans \emph{four} consecutive
  64-bit cells; the instruction touches $3\times4=12$ memory cells (was 6).
  Operand $o$ names cells $\mathbf{fp}\cdot o,\ \gen\cdot\mathbf{fp}\cdot o,\
  \gen^2\cdot\mathbf{fp}\cdot o,\ \gen^3\cdot\mathbf{fp}\cdot o$ (three
  virtual $\times\gen$ successors). Bus traffic per hash row grows
  accordingly; the value cells are still not re-committed but read out of the
  flock witness (\S\ref{sec:pipeline}).
\item Public input: $256$ bits now occupy memory cells $0..3$ (a 2-variable
  subcube); the verifier's line $\widetilde{\mathrm{PI}}$ becomes a
  square.
\item Compiler/programs: immediates are 64-bit; XMSS digests are 4 words;
  frame layouts, trampolines, range checks are otherwise untouched.
\end{itemize}

\subsection{Protocol randomness moves to $\E$}

All challenges are sampled from $\E$ (16-byte transcript reads, as today,
in the tower representation); all sumchecks, zerochecks and GKR passes run
over $\E$. Committed columns are $\K$-valued, so:
\begin{itemize}
\item the bus fingerprint $\pi_\alpha(\sigma)=\sum_i\alpha^i\sigma_i$ has
  $\sigma_i\in\K\subset\E$, $\alpha\in\E$: Schwartz--Zippel over
  $(\alpha,\gamma)$ errs at $K(m-1)/2^{128}$, as today;
\item GKR leaves $\gamma-\pi_\alpha(\sigma)$ and all product-tree layers
  are $\E$-valued, computed with the deferred-reduction accumulators;
\item zerocheck first rounds evaluate $\K$-valued columns against $\E$-valued
  $\eq$ tensors: mixed products at a third of today's cost.
\end{itemize}

\section{One commitment over $\K$, one ring switching}\label{sec:rs}

\subsection{The stacked witness}\label{sec:lists}

As today (doc \S3.1), no column is committed on its own: all committed
data is stacked largest-first into \emph{one} multilinear
$q:\cube{M}\to\K$, zero-padded to a power of two, with each column at a
$2^\kappa$-aligned offset recoverable through an $\eq$-selector. The only
changes are the value type ($\K$ instead of $\F_{2^{128}}$) and one
region: flock's Boolean witness enters as the \emph{packed} column
$q_\pkd$, $64=2^6$ bits per $\K$-word (\texttt{LOG\_PACKING} $7\to6$),
exactly as it enters today's $\F_{2^{128}}$ stack at $128$ bits per word.

\subsection{Base-field commitment, extension-field opening}\label{sec:pipeline}

Ligerito commits $q$ \emph{over $\K$} and opens it \emph{over $\E$}: an
opening proves $\sum_{w\in\cube{M}} W(w)\,\widehat q(w)=C$ for any
verifier-evaluable weight $W:\cube{M}\to\E$. This base-field-commit,
extension-field-open mode is standard in this PCS family; WHIR~\cite{whir}
(the very similar PCS used by leanVM) ships it natively. Nothing in
Ligerito resists it: the commitment layer (the Reed--Solomon code over
$\K$, the Merkle tree over 8-byte $\K$-symbols) only contributes
\emph{distance}, while every random ingredient (row batching, folding
challenges, the opening sumcheck) is sampled from $\E$ with error terms
$c/2^{128}$, exactly as today. After the first fold with an $\E$-challenge
the working vectors are $\E$-valued regardless; only the level-0 codeword
and its Merkle leaves are $\K$-valued, and query answers shrink
accordingly.

The claim pipeline is then \emph{today's, verbatim} (doc \S10): every
reduction deposits evaluation claims on committed columns into the claim
pool; each claim is an $\E$-weighted sum over the $\K$-stack via the
stacking selectors; a $\lambda$-batch folds the pool into a single
weighted claim, discharged by one batch-mixed Ligerito opening. VM columns
need no ring switching at all. The single addition, as today, is flock's
ring switching, whose output claim on the $q_\pkd$ region joins the same
opening.

\subsection{Flock's ring switching: $\F_2\to\K$, opened over $\E$}\label{sec:multi}

The bit-level claims on flock's witness are arranged to be a
\emph{single} global claim $\widehat q_F(\zeta)=v$ over $\E$ (flock's
zerocheck output; any further bit-witness claims, e.g.\ the BLAKE3 value
cells, $\K$-words read out of the bit witness by 64-bit monomial-weighted
selection, are batched into it upstream, inside flock's own reduction).
Here the note's~\cite{rsgen} two degrees genuinely differ: the packing
field is $\K$ ($f=64$), the opening field is $\E$ ($e=128$), and the note
is stated for exactly this unrelated-degrees setting.

Split $\zeta$ at the packing boundary: with members $P_j$, $j\in\cube{6}$,
\[
  v\;=\;\sum_{j\in\cube{6}}
  \eq\bigl(j,\zeta_{\mathrm{hi}}\bigr)\,
  \widehat P_{j}(\zeta_{\mathrm{lo}}) .
\]
The prover reveals the $64$ member values
$\alpha_{j}=\widehat P_{j}(\zeta_{\mathrm{lo}})\in\E$ (the
$\widehat s$-vector, now $1$\,KB), and the verifier asserts the
combination. The note then runs as written: the verifier computes the
$e=128$ coordinate targets $t_w=\sum_j\alpha_{w,j}\,b^F_j\in\K$ itself
(the bijection; nothing more is sent), batches them with powers of
$\gamma$, and one degree-2 sumcheck over $m_F$ variables against the
committed $q_\pkd$ reduces everything to a single claim
$\widehat q_\pkd(r')=\alpha'$ at $r'\in\E^{m_F}$, which joins the pool.
Two tower-specific points:
\begin{itemize}
\item the batching challenge must come from an extension of the
  \emph{packing} field (note, Remark~1): $\gamma\leftarrow\E\supseteq\K$,
  which the tower provides by construction;
\item the sumcheck summand pairs the $\K$-valued $q_\pkd$ with the
  $\E$-valued transparent factor $B$: mixed products, $2$ \texttt{PMULL}
  per term.
\end{itemize}
Error: $(2m_F+128)/2^{128}$. The verifier's $\widetilde B(r')$ is one
matrix-branching-program evaluation of width $e=128$, the DP24
tensor-algebra machinery already in \texttt{pcs/tensor\_algebra.rs}.
Revealing the member values is safe for one reason only: they feed the
bijection, so a fake vector forces some $t_w$ wrong with certainty and the
$\gamma$-batched sumcheck rejects; the values must never be checked only
through a combined opening. (The prototype in
\texttt{tests/ring\_switch\_proto.rs} exercises exactly this
reveal-and-assert machinery and its attacks, end to end.)

\section{Ligerito over $\K$}\label{sec:pcs}

\paragraph{Code layer.} One additive NTT over $\K$, encoding
$\K$-messages into $\K$-codewords: single lane, no interleaving. (New
module, the $\F_{2^{64}}$ analogue of \path{ntt/additive_ntt_f128.rs};
LCH basis over an $\F_2$-subspace of dimension $\le 64$: domains up to
$2^{64}$, far beyond need.) Symbols shrink to $8$ bytes; the Merkle layer is otherwise
untouched. Folds pair $\K$-symbols with $\E$-challenges (mixed kernels)
at the first level and are $\E$-valued afterwards, where deeper Ligerito
levels commit $\E$-valued codewords (two $\K$-lanes, encoded lane-wise by
the same NTT).

\paragraph{Parameters.} Unchanged. The challenge field is the same
$2^{128}$, the rate and decoding regime carry over, and so do query counts
and proof-of-work grinding; the base-field commitment only shrinks the
level-0 query answers. The base-commit analysis (correlated agreement for
$\E$-linear combinations of $\K$-codewords) is the standard one this PCS
family already relies on, stated explicitly by WHIR~\cite{whir}; \WP4
writes it down for our parameters.

\section{Soundness budget}\label{sec:soundness}

All interactive reductions sample from $\E$; writing $s$ for constraint
counts, $\tau$ for table log-sizes, $K$ for bus tuples, $\mu$ for GKR tree
height, $J$ for the pool size, every term keeps its current form
$c/2^{128}$: constraint batching $(s-1)/2^{128}$; zerochecks
$O(\tau)/2^{128}$; bus fingerprint $K(m-1)/2^{128}$; GKR
$O(\mu^2)/2^{128}$; flock's ring switching $(2m_F+128)/2^{128}$
(\S\ref{sec:multi}); the pool batch $J/2^{128}$. The counting arguments
(memory checking, count non-wrap, range checks) contribute no
probabilistic term at all: under the instance caps of
\S\ref{sec:vm-transport} they are exact. The PCS keeps today's
parameters. \textbf{End-to-end: $\sim$120 bits, unchanged}: this
transition moves the data and commitment field, not the security level;
the upgrade past $128$ bits is the degree-3 outlook.

\section{Performance forecast}\label{sec:perf}

Directional accounting against today's prover, in \texttt{PMULL} counts
(\S\ref{sec:fields}):

\begin{itemize}
\item \textbf{Committed data halves.} Words shrink $128\to64$ bits; the
  flock side was already bit-packed. Half the bytes to encode and hash;
  8-byte symbols halve the level-0 query answers.
\item \textbf{Trace generation and witness building} run in $\K$: one
  \texttt{PMULL} per product, $6\times$ cheaper than today.
\item \textbf{Sumcheck/zerocheck early rounds} pair $\K$-columns with
  $\E$-tensors: $2$ \texttt{PMULL} per term, a third of today. Later
  rounds and \textbf{GKR} are pure $\E$: on par with today per element.
  GKR leaf count is unchanged; BLAKE3 rows add bus tuples ($12$ vs.\ $6$
  memory flushes).
\item \textbf{NTT}: $\K$-butterflies at $1$ \texttt{PMULL} on half the
  data.
\item \textbf{No new sumchecks.} The pipeline has exactly today's shape:
  one flock ring-switching sumcheck (now with mixed $2$-\texttt{PMULL}
  terms) and one batch-mixed opening; the $\widehat s$-vector halves to
  $1$\,KB.
\end{itemize}
Net expectation: prover faster across the board, proof slightly smaller,
security level unchanged.

\section{Actionable migration plan}\label{sec:plan}

Work packages in dependency order. Each lists its primary files and an
acceptance gate. $\K$ denotes the new first-class $\F_{2^{64}}$ type
(\texttt{F64}); the tower field type is \texttt{F128T}.

\paragraph{\WP{1}: Field plumbing.}
Promote $\K$ to a first-class field type \texttt{F64} in
\path{vendor/flock-core/src/field/} (NEON mul $=1$ \texttt{PMULL} $+$
fold, \texttt{square}, \texttt{inv}, serde). Implement the tower type
\texttt{F128T} $=\K[y]/(y^2+y+x^{61})$: 2-term Karatsuba mul,
deferred-reduction accumulator, \texttt{square}, \texttt{inv}, mixed ops
\texttt{mul\_base(F64)} at $2$ \texttt{PMULL}, \texttt{From<F64>}, batched
kernels. Unit tests mirror \texttt{gf2\_64x3.rs} (portable references,
independently generated vectors, primitivity of $x$, the trace-1 check on
$c$). Add \texttt{F128T} rows to \texttt{field\_bench} against the GHASH
\texttt{F128}. \texttt{F192} stays in the tree for the outlook.
\emph{Gate:} field tests green, including mixed-op agreement with the
portable references; \texttt{F128T} multiplication on par with GHASH.

\paragraph{\WP{2}: Transcript and challenger.}
\path{src/transcript.rs} and \path{flock-core}'s \path{challenger.rs}:
sample $\E$ challenges (16 bytes, tower representation), domain
separation, keep BLAKE3 hashing and PoW grinding.
\emph{Gate:} transcript vectors regenerate deterministically; grinding test.

\paragraph{\WP{3}: Additive NTT over $\K$ and $\K$-symbol commitment.}
New \path{ntt/additive_ntt_f64.rs} (LCH over $\K$, parallel version);
$\K$-codeword commitment with 8-byte symbols (\path{pcs/commit.rs},
\path{merkle.rs}).
\emph{Gate:} encode/decode round-trip tests; distance spot-checks against a
naive Lagrange evaluator; NTT bench.

\paragraph{\WP{4}: Ligerito, $\K$-commit / $\E$-open.}
\texttt{pcs/ligerito.rs}, \texttt{pcs/basefold.rs}: level-0 fold kernels
pair $\K$-symbols with $\E$-challenges (mixed products, deferred
reduction); deeper levels are $\E$-valued (two $\K$-lanes through the same
NTT). Keep today's parameter derivation, query counts and grinding bits;
write down the base-commit correlated-agreement statement for our
parameters (WHIR~\cite{whir} precedent); \texttt{src/pcs.rs}
\texttt{PROFILE}/\texttt{MIN\_MU} untouched.
\emph{Gate:} commit/open/verify round-trip at several sizes; opening bench
against today's numbers.

\paragraph{\WP{5}: VM in $\K$.}
The \texttt{src/cpu} modules (\texttt{isa}, \texttt{execute},
\texttt{trace}, \texttt{layout}): 64-bit machine words, $\gen=x\in\K$,
BLAKE3 operands as 4-cell values (12 reads);
\texttt{src/field.rs}: \texttt{g\_powers}, \texttt{x\_pow},
\texttt{index\_mle} over $\K$ (evaluated into $\E$);
\texttt{src/compiler}: 64-bit immediates, 4-word digests;
\texttt{src/tables.rs} and \texttt{src/constraints.rs}: $\K$-valued
columns, $\E$-valued zerocheck (mixed first rounds); \texttt{src/gkr.rs}
and \texttt{src/leaf.rs}: $\E$-valued trees over $\K$-fingerprint inputs;
\texttt{src/witness.rs}: the $\K$-stack; \texttt{tests}: xmss\_vm,
hash\_chain, range\_check ported. The verifier gains the
\emph{instance-cap checks} of \S\ref{sec:vm-transport}: announced table
sizes, memory size $h$, and program length $N$ are validated against the
caps before any reduction runs.
\emph{Gate:} full VM test suite green with commitments stubbed (claims
checked directly against columns); negative tests for each cap (an
announcement exceeding $h$, $N$, or the row cap must be rejected).

\paragraph{\WP{6}: Flock's ring switching, $\F_2\to\K$.}
Adapt the DP24 pipeline (\path{pcs/ring_switch.rs}, \path{pcs/pack.rs},
\path{pcs/tensor_algebra.rs}, \path{flock-prover}'s witness layout) from
$f=e=128$ to $f=64$, $e=128$: \texttt{LOG\_PACKING} $7\to6$, the
$\widehat s$-vector at $64$ entries, the coordinate targets $t_w\in\K$,
the batching challenge $\gamma\in\E$ (note, Remark~1), the sumcheck
summand mixed $\K\times\E$; the univariate-skip prefix accounting moves
from $7$ to $6$ packing variables, and every representation-sensitive
constant (the $\varphi_8$ LCH basis, the \texttt{rs\_eq} machinery) is
re-derived for $\K$ and the tower. Ensure every bit-witness claim,
including the BLAKE3 value cells, is batched into flock's single output
claim upstream. The prototype in \texttt{tests/ring\_switch\_proto.rs}
provides reference machinery for the reveal-and-assert step and its
adversarial tests.
\emph{Gate:} flock's BLAKE3 R1CS proof verifies end-to-end against the
$\K$-packed witness; cross-check the packed polynomial against a bit-level
reference.

\paragraph{\WP{7}: Protocol assembly.}
In \texttt{src} (\texttt{pcs.rs}, \texttt{blake3\_flock.rs},
\texttt{lib.rs}, \texttt{main.rs}): the single $\K$-stack (VM columns plus
$q_\pkd$), the pool and $\lambda$-batch as today, flock's ring-switched
claim routed into the same batch-mixed opening, proof types/serde
updated, proof sizes re-measured.
\emph{Gate:} \texttt{cargo testall} green (xmss\_vm, hash\_chain
end-to-end); proof-size and prover-time table added to README; soundness
budget of \S\ref{sec:soundness} re-checked against final parameters.

\paragraph{Suggested sequencing.} \WP1$\to$\WP2 unblock everything;
\WP3$\to$\WP4 (PCS) and \WP5 (VM) proceed in parallel; \WP6 needs
\WP1--\WP4; \WP7 last. The system stays green throughout if \WP5's gate
runs with stubbed commitments, letting the VM port land before the PCS
port.

\section{Risks and open questions}\label{sec:risks}

\begin{itemize}
\item \textbf{Base-commit analysis} (\WP4). Committing over $\K$ and
  opening over $\E$ is standard (WHIR~\cite{whir} ships it), but the
  correlated-agreement statement for $\E$-linear combinations of
  $\K$-codewords must be written down for our concrete Ligerito
  parameters, not assumed; have it reviewed before freezing.
\item \textbf{Representation swap breadth} (\WP2, \WP4, \WP6). Everything
  typed \texttt{F128} changes meaning: the GHASH and tower representations
  must be distinct types with no byte-level interchange, and every
  representation-sensitive constant (LCH/$\varphi_8$ bases,
  \texttt{rs\_eq} tables, transcript encoding) re-derived. Mechanical, but
  wide; the type system should do the policing.
\item \textbf{$\E\times\E$ parity} (\WP1). The tower product is expected
  on par with GHASH ($3$ Karatsuba \texttt{PMULL} plus folds versus
  GHASH's $6$); confirm on hardware with \texttt{field\_bench} before
  committing to the swap everywhere.
\item \textbf{Ring-switching rework depth} (\WP6). Going from $f=e=128$ to
  $f=64$, $e=128$ changes the type of the coordinate targets ($t_w\in\K$)
  and the packing prefix width; the univariate-skip coupling in
  \path{pcs/ring_switch.rs} is the place where surprises would live.
\item \textbf{Deferred security upgrade.} This transition keeps
  $\sim$120 bits. The degree-3 path to more than $128$ bits is implemented
  (\texttt{F192}) and its non-power-of-two ring switching is covered by
  the note~\cite{rsgen}; revisit once the 64-bit data plane is landed.
\item \textbf{Non-goals for this transition}: no change to BLAKE3 as the
  hash, to the M3/bus design, to the ISA surface beyond the word width,
  to the PCS parameters, or to the XMSS scheme.
\end{itemize}

\begin{thebibliography}{9}
\bibitem{rsgen} Ring switching, generalized. This repository,
\texttt{ring-switching-generalized.tex}.
\bibitem{DP24} Benjamin E. Diamond and Jim Posen. Polylogarithmic Proofs for
Multilinears over Binary Towers. Cryptology ePrint Archive, Paper 2024/504.
\url{https://eprint.iacr.org/2024/504}
\bibitem{ligerito} Andrija Novakovic and Guillermo Angeris. Ligerito: A Small
and Concretely Fast Polynomial Commitment Scheme. Cryptology ePrint Archive,
Paper 2025/1187. \url{https://eprint.iacr.org/2025/1187}
\bibitem{whir} Gal Arnon, Alessandro Chiesa, Giacomo Fenzi, and Eylon
Yogev. WHIR: Reed--Solomon Proximity Testing with Super-Fast Verification.
Cryptology ePrint Archive, Paper 2024/1586.
\url{https://eprint.iacr.org/2024/1586}
\bibitem{flock} Benedikt B\"unz, Ron Rothblum, and William Wang. Flock: Fast
proving for batch Boolean computations. Cryptology ePrint Archive, Paper
2026/1329. \url{https://eprint.iacr.org/2026/1329}
\bibitem{doc} leanVM-b specification. This repository,
\texttt{misc/doc.tex}.
\end{thebibliography}

\end{document}
